\section{What happens to the continuum?}
Every real algebraic point is a genuine Euclidean point. The embedding
$K^2\hookrightarrow\mathbb R^2$ preserves and reflects \emph{every}
first-order point formula on its parameters, by Theorem~\ref{thm:transfer}.
This property is called an elementary embedding. In the standard model
$K$ is itself Archimedean: read the coarse enclosure of $x$ and take an
integer exceeding its endpoints in absolute value. For $x<y$ the exact
sign algorithm provides a positive rational lower bound on the gap;
sufficiently fine rational reads give disjoint intervals with a rational
between them. This proves rational density constructively.

Completion nevertheless changes the individual points. Externally, finite
codes are countable, whereas the points of a real line are uncountable.
For a concrete second-order distinction take a transcendental real $t$.
The external subset $\{a\in K:a<t\}$ has no supremum in $K$. A proposed
supremum $s<t$ is too small by rational density, and for $s>t$ there
is a smaller rational upper bound. The remaining equality would make
$t$ algebraic. This subset cannot be first-order definable over $K$ with
algebraic parameters, since the one-variable cell argument gives a
supremum for each inhabited bounded definable subset.

\paragraph{A separate presentation for approximations.}
One can work over a real coordinate specified by a particular PR family
of nested rational intervals $[l_n,r_n]$ of width at most $2^{-n}$.
Such specifications may carry free numerical parameters. Coordinate
equality means agreement at every precision. A strict sign has a
positive separation witness. In this setting affine combinations, dot
products, projections with a supplied lower bound on $N(u)$, and
well-separated line intersections retain effective meanings. For example,
for $d=[u,v]$ with $|d|\ge\eta>0$, an enclosing interval for $d$
of width $<\eta/2$ has distance at least $\eta/2$ from zero. Reciprocal
then magnifies interval widths by at most $4/\eta^2$. Applying the
sum and product width bounds of Section~\ref{sec:scalar} to
$A+u[B-A,v]/d$ gives a PR request for input precisions at any desired
output precision. Each fixed finite construction is handled by passing
accuracy requests backwards along its operation list. At exact tangency
the intersection point itself may still be given; a uniform test for the
number of intersections has a different computational requirement.

For a proof, fix an effective enumeration of machines on empty input.
Bounded simulation for $n$ steps is primitive recursive in the machine
index and $n$. If the first halt occurs at step $t\le n$, let
$I_{e,n}=[2^{-t},2^{-t}]$. Otherwise let $I_{e,n}=[0,2^{-n}]$.
These are nested, have the required width, and represent a scalar
$x_e$ equal to zero exactly when machine $e$ never halts. A uniform
effective test for equality of all such names would decide the halting
problem. Briefly, a halting decider on machines and inputs would allow a
machine on input $k$ to halt just when machine $k$ on input $k$ is
predicted to run forever; applying this to its own index contradicts
either prediction. Compiling any fixed input into an empty-input
machine transfers that impossibility to the present enumeration.
All the individual $x_e$ are even rational in the classical metatheory.
What fails is uniform conversion of these names to exact algebraic comparison
data. Geometrically, the unit circle and horizontal line $y=1+x_e$ are
tangent exactly in the zero case and disjoint otherwise. This shows why
the choice of input representation matters even to an elementary
construction. See Beeson's analysis of uniform geometric constructions
for the different role of continuity in synthetic constructive geometry
\cite{beeson}.
